\scanpage{284}
\mspage{14}

\deleted{Soit} Prouvons la surjectivité de $\operatorname{Hom}_M(e,X)\unsure{{-}|\longrightarrow}\operatorname{Hom}_W(e,X)$ pour $X\ne\varnothing_M$, en montrant que tout chemin réduit de longueur $2n+1$ est ``homotope'' à un chemin réduit de longueur $2n-1$ :

\diagram{\begin{tikzcd}[ampersand replacement=\&,column sep=small]
e=X_0 \arrow[r,"{\deleted{\unsure{\leftarrow}}}" description] \& X_{\unsure{1|p}} \& X_2 \arrow[l,"W"'," {\deleted{\unsure{\rightarrow}}}" description] \arrow[r,dash] \& X_3 \& \cdots \arrow[l] \& X_{2n} \arrow[l,"{\deleted{\unsure{\rightarrow}}}" description] \arrow[r,"{\deleted{\unsure{\leftarrow}}}" description] \& X_{2n+1}=X
\end{tikzcd}}

En effet, on peut remplacer la section

\diagram{\begin{tikzcd}[ampersand replacement=\&,column sep=small]
e=X_0 \arrow[r,"{\deleted{\unsure{W}}}","u_1"'," {\deleted{\unsure{\leftarrow}}}" description] \arrow[rr,bend right=40,dashed,"s"'] \arrow[r,bend right=22,"{\deleted{\unsure{s}}}"'] \& X_1 \& X_2 \arrow[l,"W"',"u_2"] \arrow[r,"{\deleted{\unsure{W}}}","u_3"'] \& X_3 \& X_4 \arrow[l,"W"']
\end{tikzcd}}

\deleted{\unsure{puis} d'abord par}

\deleted{\diagram{\begin{tikzcd}[ampersand replacement=\&]
\unsure{e} \arrow[r,"{\unsure{u_2s}}"] \& \unsure{X_2} \& \unsure{X_3}\arrow[l,"{\unsure{W}}"'] \& \cdots
\end{tikzcd}}}

par

\diagram{\begin{tikzcd}[ampersand replacement=\&]
e\arrow[r,"{\deleted{\unsure{u_3s}}}","{u_3s}"'] \& X_3 \& \cdots \arrow[l,"W"']
\end{tikzcd}}

\deleted{(où $s$ est une section de $X_1$ sur $e=X_0$), puis par}

où $s$ est une section de $X_2$ sur $e$ telle que

\[
u_2s\deleted{\unsure{\simeq}}=u_1\qquad\text{\uuline{dans }}MW^{-1}.
\]

Cela prouve donc par \unsure{réc} sur $n$ que le chemin est homotope à un chemin de longueur $1$, $e\longrightarrow X$, qed.

Il faut donc vérifier, \deleted{1\textsuperscript{o}} dans le cas actuel, vérifier \uline{1\textsuperscript{o}}) (c'est le \inserted{cor.} \unsure{cor.} \unsure{2|?}.12.), et 2\textsuperscript{o})

\diagram{\begin{tikzcd}[ampersand replacement=\&]
X_1 \& X_2 \arrow[l,"{u_2\in W}"']\\
e \arrow[u,"u_1"] \arrow[ur,dotted,"s"'] \&
\end{tikzcd}}

Mais comme $u_2\in W$, on a $\pi_0(u_2)$ \unsure{iso|bij.}, donc

% archive sheet number: 283
